\chapter{系统详细设计与实现}

\section{开发环境与工具}

本文系统在 Windows 环境下完成开发、联调与论文撰写。后端采用 Java 17 与 Spring Boot 2.7 构建，前端采用 Vue 3 与 Vite 进行单页应用开发，数据库使用 MySQL 持久化业务数据，项目通过 Maven 与 Node.js 管理依赖和构建流程。除基础开发环境外，系统还借助 JUnit 5、Mockito、浏览器开发者工具以及前端联调脚本完成测试与验证。系统主要开发环境与工具如表~\ref{tab:dev-env-real} 所示。

\begin{table}[htbp]
\centering
\footnotesize
\caption{系统开发环境与主要工具}
\label{tab:dev-env-real}
\begin{tabular}{|c|c|p{6.2cm}|}
\hline
类别 & 名称/版本 & 作用说明 \\\hline
操作系统 & Windows & 系统开发、调试、论文编写与本地集成测试环境。 \\\hline
后端语言 & Java 17 & 实现认证授权、业务服务、数据处理与接口层逻辑。 \\\hline
后端框架 & Spring Boot 2.7.13 & 构建 RESTful 服务并整合安全、数据访问与消息组件。 \\\hline
安全框架 & Spring Security & 完成接口访问控制、身份恢复与角色限制。 \\\hline
令牌机制 & JWT & 在前后端分离架构下实现无状态认证。 \\\hline
持久化技术 & Spring Data JPA & 实现实体映射、仓储接口与事务访问。 \\\hline
数据库 & MySQL & 存储订单、库存、采购、生产、质检和计划等业务数据。 \\\hline
前端框架 & Vue 3 & 构建多角色业务页面与组件交互。 \\\hline
前端构建工具 & Vite & 提供快速开发启动、热更新与生产构建能力。 \\\hline
HTTP 通信 & Axios & 统一封装前后端接口调用、令牌携带和异常处理。 \\\hline
实时通信 & SockJS、STOMP、WebSocket & 实现消息主题订阅与事件实时推送。 \\\hline
测试工具 & JUnit 5、Mockito、Node.js 脚本 & 完成后端单元测试与前端接口/联调验证。 \\\hline
\end{tabular}
\end{table}

从工程组织角度看，后端源码主要位于 \texttt{src/main/java/com/code} 目录，按照 \texttt{controller}、\texttt{service}、\texttt{repository}、\texttt{entity} 和 \texttt{security} 等子包进行划分；前端源码位于 \texttt{frontend/src}，按照 \texttt{pages}、\texttt{api}、\texttt{router} 和 \texttt{store} 等方式组织。这种目录结构为系统模块扩展与后续维护提供了清晰基础。

\section{各功能模块实现}

为使论文内容更贴近项目实际，本节不再停留于概括性描述，而是从“前端页面入口—后端控制器—业务服务—关键实体”的映射关系出发，对系统主要功能模块进行说明。系统主要模块及其实现入口如表~\ref{tab:module-file-mapping} 所示。

\begin{table}[htbp]
\centering
\footnotesize
\caption{主要功能模块与实现入口对应关系}
\label{tab:module-file-mapping}
\begin{tabular}{|c|p{3.2cm}|p{3.3cm}|p{4.3cm}|}
\hline
模块 & 典型前端页面 & 主要控制器 & 核心服务/说明 \\\hline
认证与账号管理 & Login、Register、AccountAdmin & AuthController、AdminController、ProfileController & SecurityConfig、JWT 认证、账号维护逻辑 \\\hline
订单与销售 & Customer、Orders、SalesRecords & CustomerController、SalesOrderController & OrderWorkflowService、销售记录生成 \\\hline
库存与预警 & Inventory、InventoryAlert & InventoryController、InventoryAlertController & 库存出入库、预警分析与联动 \\\hline
生产管理 & Production、ProductionPlan & ProductionController & OrderWorkflowService、WeeklyPlanningService \\\hline
采购与供应商 & Procurement、ProcurementPlan、RawMaterials & ProcurementController & ProcurementWorkflowService、周采购计划 \\\hline
质检与追溯 & Quality、Customer & QualityController、CustomerController & QualityService、批次追溯与异常回流 \\\hline
实时通知 & Dashboard、各角色待办区 & 多控制器共享消息入口 & WebSocketConfig、NotificationService \\\hline
\end{tabular}
\end{table}

\subsection{用户认证与账号管理模块实现}

用户认证模块由后端安全配置、认证接口和前端登录态管理共同构成。后端在 \texttt{SecurityConfig} 中定义安全过滤链，并对 \texttt{/api/v1/auth/**}、\texttt{/api/v1/admin/**}、\texttt{/api/v1/customer/**}、\texttt{/api/v1/orders/**}、\texttt{/api/v1/procurement/**} 等不同路径进行角色访问限制。用户提交登录请求后，认证管理器根据账号和密码完成校验，校验成功后生成带有用户身份与角色信息的 JWT 并返回前端。前端在登录成功后保存令牌，在后续通过 Axios 访问接口时统一附加 \texttt{Authorization: Bearer <token>} 请求头。

在项目实现中，认证模块具有以下特点。第一，系统采用无状态认证方式，后端不为用户维护传统服务端 Session，而是通过令牌恢复当前用户身份。第二，角色控制不仅体现在页面菜单层面，更体现在接口访问层面。第三，系统支持系统管理员维护内部账号，能够查看当前内部人员列表、按角色筛选账号、创建新账号、修改角色与启停状态。该功能由 \texttt{AdminController} 实现，其数据对象直接来源于 \texttt{User} 和 \texttt{Role} 实体，并通过后端规则保证系统中至少保留一个启用状态的管理员账号。

此外，账号管理与认证模块还承担统一术语与组织结构映射职责。系统内部将角色统一保存为 \texttt{ROLE\_*} 形式，而在前端页面中再映射为“销售管理员”“仓库管理员”“采购管理员”等中文标签，以兼顾程序内部标准化与用户界面可读性。由此可见，认证授权模块既是安全入口，也是系统角色化设计的基础。

\subsection{顾客下单与订单管理模块实现}

顾客下单功能以顾客页面和订单接口为核心实现。前端下单时，用户选择产品、填写数量和收货信息，页面在本地完成基础校验后调用订单创建接口。后端在 \texttt{SalesOrderController} 中接收订单数据，对订单明细逐项计算行金额与总金额，并在保存订单主表和明细表时建立 \texttt{SalesOrder} 与 \texttt{OrderItem} 的从属关系。订单创建成功后，系统立即将订单写入数据库，并通过 WebSocket 向订单主题发送新订单消息。

订单管理并非单一角色独占功能，而是顾客、销售管理员、仓库管理员和生产管理员共同参与的协同流程。因此，订单实体 \texttt{SalesOrder} 的核心不只是记录产品和金额，还需要承载流程状态。项目当前实现中，订单常见状态包括“待销售审核”“待仓库核查”“已接单”“生产中”“待质检”“待生产入库”“已发货”等。所有状态变更均由后端服务 \texttt{OrderWorkflowService} 统一控制，避免前端直接修改数据库中的状态值。

对于顾客而言，系统除下单外还支持查看本人订单。当前项目已实现顾客查看自己下单记录的功能，展示内容包括产品名称、数量、金额、订单编号和状态等。这一功能增强了顾客侧信息透明度，也为后续质量追溯提供了入口。

\subsection{销售管理模块实现}

销售管理模块主要承担订单审核与销售记录沉淀功能。销售管理员在订单列表中可查看顾客提交的待处理订单，并基于订单信息执行接受或拒绝操作。当销售接受订单时，控制器会调用 \texttt{OrderWorkflowService.routeToWarehouseCheck} 将订单推进为“待仓库核查”；当销售拒绝订单时，系统要求填写拒绝理由，并将订单状态更新为“已拒绝”。这一设计体现了订单入口把关机制，避免不合理订单直接进入库存和生产流程。

销售记录功能则用于沉淀已完成销售业务的结果信息。项目中，当订单最终进入“已完成”状态后，系统在 \texttt{SalesOrderController} 内调用 \texttt{createSalesRecordIfNeeded} 自动生成 \texttt{SalesRecord} 记录。该记录包含销售记录编号、订单号、顾客名称、地址、订单金额、状态以及操作销售管理员的编号和姓名等信息。为满足岗位隔离要求，普通销售管理员仅查看自己创建的销售记录，而系统管理员可以查看全量销售记录。系统同时支持销售记录的时间范围筛选和 CSV 导出，说明该模块不仅服务于业务操作，也服务于事后统计分析。

从实现效果看，销售模块的关键价值在于将“订单入口审核”和“销售结果留痕”结合起来。一方面，订单审核决定是否进入后续仓储执行；另一方面，销售记录保留了订单完成后的业务归属与操作责任，从而为系统的过程审计与绩效统计提供数据基础。

\subsection{仓库管理与库存预警模块实现}

仓库管理模块是订单履约和采购入库的核心中间层，负责库存展示、库存检索、成品与原材料出入库、发货处理、库存流水生成以及预警分析。项目中，库存实体由 \texttt{InventoryItem} 表示，记录产品、仓库、数量、预留数量和批次号等信息；库存变化则通过 \texttt{StockTransaction} 持久化记录，保存事务编号、出入库方向、关联业务类型、关联业务编号、操作人和时间等内容。

在功能实现上，仓库模块主要包含三类处理逻辑。第一类是订单驱动的库存核查与发货。仓库接到已审核订单后，系统通过订单工作流服务计算产品可用库存，即 \texttt{quantity - reservedQuantity}。若库存足够，则执行预留库存；若库存不足，则生成生产计划。发货阶段系统优先扣减已预留库存，若预留不足再消耗剩余可用库存，并同步写入出库流水。第二类是采购或生产驱动的入库。采购到货时，仓库确认收货后系统为原材料增加库存并生成入库流水；生产批次质检合格后，仓库执行生产入库确认，将计划转为已入库状态。第三类是库存查询与检索。系统支持按成品和原材料维度展示库存信息，并显示库存名称、数量、可用量、仓库信息与预警状态。

库存预警由 \texttt{InventoryAlertController} 提供支持。该模块分别对成品和原材料进行短缺分析：对成品，比较安全库存与当前可用库存，并考虑在制生产量；对原材料，则在安全库存与当前可用库存基础上继续考虑在途采购量。当短缺数量大于零时，系统生成预警项，并为仓库管理员提供“创建生产计划”或“创建采购申请”的推荐动作。由此可见，库存预警模块将仓库管理从事后登记提升到了事前感知与联动处置层面。

\subsection{生产管理模块实现}

生产管理模块以订单缺口和计划建议为触发源，围绕生产计划生成、生产执行、生产完工回传和生产记录展示展开。当前项目中，订单库存不足时，\texttt{OrderWorkflowService} 会依据缺口数量创建 \texttt{ProductionPlan} 记录，并以“PLAN-订单号-产品编号-时间戳”的方式生成计划编号。计划实体保存产品、计划数量、状态、开始时间、结束时间、创建人及完成人等信息，这使系统能够明确每条计划的来源与责任归属。

生产管理员在前端页面中查看计划后，可执行生产完工操作。完工并不意味着流程结束，而是意味着系统进入“生产结束、等待质检”的新阶段。项目实现中，生产完工时系统将计划状态更新为 \texttt{DONE}，记录完工人员信息，并根据订单号和产品生成或更新 \texttt{Batch} 批次，随后广播质检待办通知。此时订单状态会更新为“待质检”，从而阻止仓库在未检验前进行入库或发货。

除了订单驱动的生产任务外，系统还支持从库存预警和生产周计划角度生成计划建议。也就是说，生产模块既响应短期订单履约缺口，也服务于按周维度的生产准备工作。这种设计增强了系统对不同时间尺度生产需求的适应能力。

\subsection{采购与供应商管理模块实现}

采购与供应商管理模块主要实现从采购单创建到供应商处理再到仓库收货入库的完整链路。采购管理员通过前端页面创建采购单时，系统会根据采购明细自动计算总金额，并通过 \texttt{ProcurementWorkflowService.createPurchaseOrder} 自动生成采购单号 \texttt{poNo}。项目中采购单号采用“PO + 时间戳 + 随机后缀”的方式生成，避免人工填写和重复冲突。

在创建采购单过程中，系统执行了较为严格的业务校验。首先，采购明细必须是产品类型为 \texttt{RAW\_MATERIAL} 的原材料；其次，原材料档案必须已经维护优选供应商信息；再次，所选供应商必须与该原材料档案中的供应商信息匹配。只有满足上述条件，采购单才能被成功创建并进入“待供应商确认”状态。该校验逻辑充分体现了项目针对“采购管理员选择的供应商应与原材料档案保持一致”的业务要求。

供应商端的实现遵循“只看与自己相关的数据”原则。供应商登录后，仅能查看与本人账号匹配的原材料档案和采购订单。供应商可对采购单执行接单、拒单或发货操作，每次状态变化都会触发对采购管理员的实时通知。采购管理员在接收到供应商发货通知后，可进一步通知仓库收货。仓库确认后，系统自动对采购明细执行入库并生成 \texttt{StockTransaction} 流水。整个采购流程的核心不是单据生成本身，而是企业内外主体围绕状态变化进行的持续协同。

\subsection{质检与质量追溯模块实现}

质检模块以 \texttt{Batch} 和 \texttt{QualityRecord} 两类实体为核心。生产完成后，系统创建与订单、产品和生产责任人相关联的批次记录，随后由质检管理员对该批次执行检验。质检结果仅支持“合格”和“不合格”两种类型，但系统要求在不合格时必须填写原因，以便后续返工和责任界定。质检操作完成后，系统同时更新批次状态、写入质检记录、补充质检员信息，并根据结果触发不同的后续流程。

若质检结果为不合格，\texttt{QualityService} 会向对应生产管理员主题发送返工通知，并将相关订单状态退回“生产中”；若质检结果为合格，系统则将订单推进为“待生产入库”，并通知仓库管理员对合格批次完成入库确认。该设计保证了质量环节对业务流程的实质性约束，而非仅作为附属登记信息存在。

质量追溯功能面向顾客侧开放。顾客可在质量追溯页面输入订单号，由系统通过订单号查询对应批次及其质检结果，并返回合格或不合格状态。由于批次记录中保存了来源订单号，质检记录又保存了批次和检验信息，因此系统能够以较低的实现复杂度完成订单级质量追溯。这种设计不仅增强了系统透明度，也体现了业务数据在设计阶段预留追溯链路的重要性。

\subsection{实时消息通知模块实现}

为了减少多角色协同中的等待时间，系统采用 WebSocket/STOMP 实现实时消息通知。后端通过 \texttt{WebSocketConfig} 注册 \texttt{/ws} 连接端点，并启用 \texttt{/topic} 与 \texttt{/queue} 主题代理；前端在用户登录后建立长连接，并根据当前角色订阅订单、生产、采购、质检、库存预警等相关主题。项目中已使用的主题包括 \texttt{/topic/orders}、\texttt{/topic/production}、\texttt{/topic/procurement}、\texttt{/topic/quality} 和 \texttt{/topic/mrp} 等。

从业务实现角度看，实时消息并不是单独存在的辅助模块，而是嵌入在各工作流服务中的基础能力。例如，销售审核通过后，系统向仓库广播“接收一条新的订单 + 订单号 + 请查看”；仓库发货后，系统向顾客广播“您的订单 + 订单号 + 已发货”；采购单创建后，系统向对应供应商广播待处理信息；生产完工后，系统向质检管理员广播待检批次；质检不合格时，系统向对应生产管理员广播返工通知。通过这种方式，系统显著减少了依赖人工转述的环节，提高了整体协同效率。

\subsection{周计划智能生成模块实现}

周计划模块由 \texttt{WeeklyPlanningService} 实现，分别负责生产周计划与采购周计划生成。与传统只针对当前事务做响应的方式不同，周计划模块尝试基于已有业务数据对未来一周的资源需求进行提前准备，因而具有一定的决策支持属性。

在生产周计划生成过程中，系统首先统计上周完工量，其次汇总当前未完成订单需求，再统计当前可用成品库存。对每个成品，系统以“上周完工量乘以增长系数 1.10”和“当前未完成订单需求”两者中的较大值作为基线，再扣减当前可用库存，最终得到建议生产量。采购周计划则以生产周计划为上游输入，通过 BOM 关系换算出原材料需求，同时结合上周采购量乘以 1.05 的增长系数、安全库存缺口、当前可用原材料库存和在途采购量，计算得到建议采购量。

为便于业务理解，系统在计划主表中记录算法说明，在计划明细中记录建议数量、上周产量或采购量、当前需求、可用库存、在途数量和建议原因。管理员和相关岗位不仅可以查看计划结果，还可以结合实际情况对建议进行人工执行或调整。这种实现方式虽然不属于复杂机器学习算法，但具有可解释性强、实现成本低和适合中小型企业业务落地等优点。

